\section{Propositional rejection and proofs of nonvalidity}
\label{sec:p02-classical-rejection}

A refutation calculus can certify that a formula fails to be a theorem
without proving its negation. This is a different task from deriving a
contradiction from assumptions. The classical construction below uses
the antiaxiom and rejection operations recorded in
\cite[\S1, pp.~68--69]{Citkin2013JankovRefutation}. The general signed
metastatement notation occurs in that paper's Definitions~4.1--4.3.

\begin{hbdefinition}{def:p02-crf-judgments}{Assertion and rejection are metalinguistic judgments}
Use the \(\{\neg,\to\}\) language and finite Hilbert derivations of
\cref{def:p02-hilbert-derivation}. Boolean valuations take values
\(0,1\), with negation exchanging these values and material implication
false exactly at antecedent one and consequent zero.
Write \(\vdash A\) when that
calculus derives \(A\) with no open assumptions. Its classical adequacy
is the empty-premise instance of
\cref{thm:p02-calculi-derivability-equivalence}.
Write \(\dashv A\) for the distinct formal rejection judgment defined
below. It is not the object formula \(\neg A\).
The source's \(A^{+}\) and \(A^{-}\) are signed metastatements:
their intended meanings are, respectively, that \(A\) belongs and does
not belong to the indicated logic. In this classical construction they
correspond to \(\vdash A\) and \(\dashv A\). The superscripts do not
extend the formula grammar.
\end{hbdefinition}

\begin{hbdefinition}{def:p02-crf-derivation}{The classical rejection calculus}
Fix one propositional variable \(p\). Its judgment \(\dashv p\) is
the antiaxiom. A rejection derivation is a finite sequence of judgments.
A positive line must be accompanied by a closed Hilbert derivation;
a negative line is the antiaxiom or follows from earlier lines by the
rejection modus-tollens or reverse-substitution operation on the sheet.
A substitution \(\sigma\) maps variables to formulas and extends
homomorphically through \(\neg,\to\). Replacement is simultaneous:
it does not apply \(\sigma\) again inside the inserted formulas.
Positive proofs, substitutions, and earlier-line references are part
of the certificate, not unrecorded oracle answers.
\end{hbdefinition}

\calculussheet{sheet:p02-classical-rejection}{Classical propositional rejection}{Closed positive CPC proofs; negative judgments certify nonvalidity; negation remains an object-language connective; no equality or open-hypothesis rejection rule.}
\hbrule{rule:p02-crf-antiaxiom}{Antiaxiom}{}{\dashv p}{p is the fixed propositional variable of the rejection calculus.}
\hbrule{rule:p02-crf-mt}{Rejection modus tollens}{\vdash(A\to B)\qquad\dashv B}{\dashv A}{The positive premise has a closed proof in the specified sound classical calculus.}
\hbrule{rule:p02-crf-rs}{Reverse substitution}{\dashv\sigma(A)}{\dashv A}{sigma is a simultaneous uniform substitution of formulas for variables.}


The positive premise of rejection modus tollens must be a theorem,
not a formula derived under hypotheses. Reverse substitution removes
a substitution from a rejected instance; it does not substitute freely
into an arbitrary rejection. The consequences of reversing either
restriction are exhibited below.

\begin{hblemma}{thm:p02-crf-substitution}{proved=annex:p02.crf.substitution}{Truth under simultaneous substitution}
For any Boolean valuation \(w\), let \(v(x)=w(\sigma(x))\) on
variables. Then \(w(\sigma(A))=v(A)\) for every formula \(A\).
Consequently a countervaluation of \(\sigma(A)\) supplies a
countervaluation of \(A\); a theorem \(A\) has only theorem
substitution instances.
\end{hblemma}

\begin{hbtheorem}{thm:p02-crf-soundness}{proved=annex:p02.crf.soundness}{Every rejection has a countervaluation}
If a rejection derivation ends in \(\dashv A\), some Boolean
valuation makes \(A\) false. In particular, no classically valid
formula is rejected, and no formula is both asserted and rejected
by these calculi.
\end{hbtheorem}

\subsection{Turning a countervaluation into a rejection certificate}
\label{sec:p02-crf-certificate}

\begin{hbdefinition}{def:p02-crf-constants}{Formula constants for the substitution construction}
Put \(T=(p\to p)\) and \(F=\neg(p\to p)\).
Every Boolean valuation gives \(T\) value one and \(F\) value zero.
Citkin's introductory construction instead uses the formulas
\((p\to p)\) and \((p\land\neg p)\). The latter has the same
constant-zero truth function as \(F\); \(F\) keeps the present
construction within the stated primitive Hilbert language.
For a countervaluation \(v\) of \(A\), define \(\sigma_v(x)=T\)
if \(v(x)=1\), and \(\sigma_v(x)=F\) if \(v(x)=0\).
The distinguished variable \(p\) need not be fresh because this
substitution is simultaneous and both inserted formulas are constant
as truth functions.
\end{hbdefinition}

\begin{hbtheorem}{thm:p02-crf-completeness}{proved=annex:p02.crf.completeness}{Every classically invalid formula has a rejection derivation}
If \(A\) has a countervaluation \(v\), then
\(\sigma_v(A)\to p\) is a classical theorem. Its closed positive
proof, the antiaxiom \(\dashv p\), rejection modus tollens, and
reverse substitution form a rejection derivation of \(A\).
Thus \(\dashv A\) is derivable exactly when \(A\) is not valid.
The proof supplies the substitution identity for every valuation and
explains the finite positive-proof obligation.
\end{hbtheorem}

\begin{hbexample}{def:p02-crf-example}{A complete rejection of an implication}
For \(A=(q\to r)\), choose \(v(q)=1,v(r)=0\).
The corresponding instance is \(\sigma_v(A)=(T\to F)\).
There is a closed positive proof of \((T\to F)\to p\).
Together with \(\dashv p\), the sheet's two rejection operations
give \(\dashv(T\to F)\) and then \(\dashv(q\to r)\).
The annex gives the actual Hilbert lines under the single temporary
assumption \(T\to F\), then discharges it by the proved deduction
theorem. The certificate does not assert \(\neg(q\to r)\): the
valuation \(q=0,r=0\) makes \(q\to r\) true.
\end{hbexample}

\begin{hbtheorem}{thm:p02-crf-polarity}{proved=annex:p02.crf.polarity}{Rejection, negation, and the three truth-function cases}
Rejection of \(A\) means \((\exists v)\,v(A)=0\), whereas assertion
of \(\neg A\) means \((\forall v)\,v(A)=0\).
A tautology is asserted and not rejected. A contradiction is rejected
and its negation is asserted. For a contingent formula, both it and
its negation are rejected, while neither is asserted. These two
rejections may use different countervaluations.
Forward substitution on arbitrary negative judgments and use of a
hypothesis-dependent positive implication in rejection modus tollens
would each allow rejection of a tautology.
\end{hbtheorem}

\begin{hbwarning}{def:p02-crf-wrong-directions}{Two unsound rule changes}
Forward negative substitution would send \(\dashv p\) to
\(\dashv(p\to p)\) by substituting \(p\to p\) for \(p\).
Also, under the hypothesis \(q\), the formula \(T\to q\) is
derivable. Combining that hypothesis-dependent derivation with
\(\dashv q\) as if it were a closed theorem would reject \(T\).
The calculus permits neither change. A rejected formula is not an
assumption that is false in every model under discussion.
\end{hbwarning}

\begin{hbmethod}{def:p02-crf-procedure}{Finding and checking an assertion or rejection certificate}
List the finitely many variables occurring in the input \(A\), and
evaluate \(A\) on all their Boolean assignments. If all results are
one, search the effectively enumerable closed Hilbert proofs for a
proof of \(A\). If some result is zero, construct \(\sigma_v\)
from that assignment, search for a closed proof of
\(\sigma_v(A)\to p\), and append the two rejection steps.
Check every positive proof against its actual axiom schemas and
modus-ponens indices; check each negative step against the literal
sheet and the displayed substitution. Finite truth-table search finds
the appropriate branch. Completeness guarantees termination of its
subsequent positive-proof search, but gives no efficient proof-size
bound.
\end{hbmethod}

\begin{hbtheorem}{thm:p02-crf-procedure}{proved=annex:p02.crf.procedure}{Termination and the meaning of the checked output}
With the countable effectively encoded propositional language and
the stated sound and complete positive calculus, the procedure
terminates on every formula. A checked positive output establishes
validity; a checked negative output establishes nonvalidity.
It does not conflate a short truth-table counterexample, a potentially
long syntactic rejection certificate, and a derivation of a negation.
\end{hbtheorem}

\begin{proofboundary}{def:p02-crf-scope}{The classical theorem does not transfer unchanged to other logics}
The completeness construction uses Boolean formula constants and
classical positive completeness. In intermediate or modal logics,
rejection may require different antiaxioms and characteristic formulas.
The signed framework alone neither proves their refutation
completeness nor turns the absence of a discovered proof into a
rejection certificate \cite[\S1 and Definitions~4.1--4.3]{Citkin2013JankovRefutation}.
\end{proofboundary}
